\chapter{Donde virá julgar os vivos e os mortos}

O sétimo artigo da Profissão de Fé completa a proclamação da glória de Cristo, afirmando que Ele ``virá julgar os vivos e os mortos''. Este artigo revela o sentido último da história e o destino final da humanidade. Cristo, que subiu aos céus e está sentado à direita do Pai, voltará no fim dos tempos para realizar o julgamento definitivo. O Catecismo da Igreja Católica, nos parágrafos 668 a 681, apresenta este mistério como o cumprimento do Reinado de Cristo, já iniciado mas ainda não manifestado em toda a sua plenitude. A esperança da vinda do Senhor anima a vida da Igreja peregrina.

\section{O Reinado de Cristo já presente (668--672)}

Cristo, depois da sua Ascensão, participa do poder e da autoridade de Deus. Ele é o Senhor do cosmos e da história. Em Cristo, a história humana e toda a criação são ordenadas e levadas à sua plenitude transcendente. ``Cristo morreu e ressuscitou para ser o Senhor dos mortos e dos vivos''. 

Como Senhor, Cristo é também a Cabeça da Igreja, que é o seu Corpo. Elevado aos céus e glorificado, Ele habita na sua Igreja através do Espírito Santo. A Redenção é a fonte da autoridade que Cristo, pelo Espírito, exerce sobre a Igreja. O Reino de Cristo já está presente em mistério sobre a terra como ``semente e início do Reino''.

Desde a Ascensão, o desígnio de Deus entrou na sua realização. Estamos já na ``última hora''. O Reino de Cristo manifesta-se já através dos sinais que acompanham a pregação da Igreja. No entanto, este Reinado ainda não está plenamente realizado. Ele está sob ataque das forças do mal, embora estas tenham sido definitivamente vencidas pela Páscoa de Cristo.

A Igreja peregrina vive entre o ``já'' e o ``ainda não''. Enquanto espera a manifestação plena do Reino, ela carrega os sinais deste mundo passageiro. Os cristãos, unidos a Cristo, gemem interiormente e esperam a revelação dos filhos de Deus. Eles rezam com insistência: ``Venha o teu Reino!'' e ``Maranatha! Vem, Senhor Jesus!''.

A realeza de Cristo é simultaneamente presente e futura. Já exerce o seu domínio através da Igreja, mas manifestar-se-á com poder e grande glória no seu regresso. Esta tensão escatológica é constitutiva da existência cristã. Os fiéis vivem no tempo presente com os olhos fixos na vinda gloriosa do Senhor.

\section{Virá julgar os vivos e os mortos (673--681)}

A vinda de Cristo como Juiz é um dos temas mais claramente proclamados nas Escrituras. Os profetas do Antigo Testamento falavam do ``Dia do Senhor''. Nos Evangelhos, o próprio Jesus descreve com pormenor o seu regresso como Juiz do mundo. Os Apóstolos deram grande relevo a esta doutrina nas suas pregações e escritos.

Antes do Juízo Final, diversos sinais precederão a vinda do Senhor. O Evangelho será pregado a todas as nações. Surgirá o Anticristo, o homem do pecado, que seduzirá muitos com prodígios enganadores e perseguirá a Igreja. Haverá perturbações extraordinárias na natureza e uma grande apostasia. Estes sinais não servem para calcular a data exacta, conhecida só pelo Pai, mas para manter os cristãos vigilantes.

O Juízo Final será universal. Todos os homens, vivos e mortos, comparecerão perante o trono de Cristo. Ele separará uns dos outros, como o pastor separa as ovelhas dos cabritos. Aos que estiverem à sua direita dirá: ``Vinde, benditos de meu Pai, recebei por herança o Reino preparado para vós desde a fundação do mundo''. Aos da esquerda: ``Apartai-vos de mim, malditos, para o fogo eterno''.

O critério do julgamento será a prática da caridade para com os mais pequenos. O próprio Cristo se identifica com os famintos, os sedentos, os estrangeiros, os nus, os doentes e os presos. ``Todas as vezes que fizestes isto a um destes meus irmãos mais pequeninos, foi a mim que o fizestes''. O Juízo Final revelará o bem e o mal escondidos no coração de cada um e manifestará o sentido de toda a história.

Cristo julgará como Deus e como Homem. O Pai entregou todo o julgamento ao Filho. Ele julgará o mundo com justiça. Os eleitos participarão de certa forma deste julgamento, especialmente os Apóstolos, que se sentarão em doze tronos para julgar as doze tribos de Israel. O Juízo Final justificará a providência divina, mostrando que tudo foi ordenado por uma sabedoria infinita e uma justiça perfeita.

A Igreja crê firmemente que Cristo virá em glória para julgar os vivos e os mortos. Este artigo de fé aparece em todos os antigos símbolos. A par do Juízo particular que cada alma recebe imediatamente após a morte, o Juízo Final manifestará publicamente a justiça de Deus e a verdade de cada vida. Será o momento em que a história humana receberá o seu veredicto definitivo.

A consciência deste julgamento final deve levar os cristãos a viverem em vigilância e em caridade activa. Ninguém sabe o dia nem a hora. Por isso, é necessário estar sempre preparado, fazendo o bem aos mais pequenos, nos quais o próprio Cristo se esconde. A esperança da vinda do Senhor não paralisa a acção, mas enche-a de sentido e de urgência missionária.
